\documentclass[../DoAn.tex]{subfiles}
\begin{document}

\section{Kết luận}

Đồ án đã xây dựng hệ thống \textit{Multi-tenant RAG chatbot hỗ trợ tư vấn bán hàng, so sánh và tham chiếu giá sản phẩm nội thất}. Hệ thống hướng tới việc phục vụ nhiều cửa hàng trên cùng một nền tảng, trong đó mỗi cửa hàng có dữ liệu sản phẩm, chatbot, knowledge base, hội thoại, lead và yêu cầu mua hàng riêng.

Kết quả đạt được của đồ án gồm các nhóm chức năng chính. Thứ nhất, hệ thống hỗ trợ mô hình multi-tenant, cho phép quản lý tenant, thành viên cửa hàng, chatbot, dữ liệu knowledge base, hội thoại, lead và yêu cầu mua hàng theo từng cửa hàng. Thứ hai, hệ thống tách backend nghiệp vụ và AI Service, trong đó backend xử lý nghiệp vụ, phân quyền và dữ liệu tenant, còn AI Service xử lý các luồng RAG và gọi Claude Sonnet để sinh câu trả lời. Thứ ba, chatbot hỗ trợ các luồng hội thoại chính như tư vấn bán hàng theo cửa hàng, so sánh sản phẩm và tham chiếu giá sản phẩm. Thứ tư, hệ thống cho phép tạo lead hoặc yêu cầu mua hàng từ hội thoại, giúp kết nối quá trình tư vấn với nghiệp vụ bán hàng của cửa hàng. Ngoài ra, hệ thống cũng hỗ trợ tích hợp Messenger, Telegram thông qua webhook và được đóng gói bằng Docker Compose để triển khai thử nghiệm.

Đóng góp chính của đồ án là kết hợp chatbot RAG với mô hình multi-tenant trong bài toán tư vấn bán hàng nội thất. Hệ thống cho phép chatbot tư vấn dựa trên dữ liệu riêng của từng cửa hàng, đồng thời hỗ trợ các chức năng gần với nhu cầu mua hàng thực tế như so sánh sản phẩm, tham chiếu giá và tạo yêu cầu mua hàng sau hội thoại.

Bên cạnh kết quả đạt được, hệ thống vẫn còn một số hạn chế. Dữ liệu sản phẩm và dữ liệu tham chiếu còn phụ thuộc vào phạm vi thu thập trong quá trình demo, chưa đầy đủ như hệ thống thương mại thực tế. Chức năng tham chiếu giá mới có ý nghĩa trong phạm vi dữ liệu của hệ thống, chưa đại diện cho toàn bộ thị trường nội thất. Chất lượng phản hồi của chatbot còn phụ thuộc vào chất lượng knowledge base, prompt và Claude Sonnet. Ngoài ra, hệ thống chưa triển khai đầy đủ các chức năng nâng cao như đánh giá tự động quy mô lớn, dashboard phân tích chuyên sâu, quản lý tồn kho, thanh toán và vận chuyển.

Nhìn chung, đồ án đã hoàn thành mục tiêu xây dựng một hệ thống chatbot tư vấn nội thất theo mô hình multi-tenant ở mức demo thực nghiệm. Hệ thống thể hiện được các luồng chính từ quản lý dữ liệu, xây dựng knowledge base, tư vấn sản phẩm, so sánh sản phẩm, tham chiếu giá, tạo yêu cầu mua hàng đến xử lý thông tin sau hội thoại.

\section{Hướng phát triển}

Trong thời gian tới, hệ thống có thể được phát triển theo một số hướng chính. Trước hết, cần mở rộng và chuẩn hóa dữ liệu sản phẩm từ nhiều cửa hàng hơn, bao gồm các thông tin như tên sản phẩm, giá, chất liệu, kích thước, phong cách, cửa hàng cung cấp và thời điểm thu thập. Dữ liệu đầy đủ và nhất quán hơn sẽ giúp chatbot tư vấn, so sánh và tham chiếu giá chính xác hơn.

Tiếp theo, cần cải thiện chất lượng truy xuất và sinh câu trả lời bằng cách kết hợp tìm kiếm từ khóa, tìm kiếm ngữ nghĩa và bộ lọc metadata như loại sản phẩm, ngân sách, chất liệu, kích thước hoặc phong cách. Hệ thống cũng cần có cơ chế kiểm tra dữ liệu trước khi trả lời để hạn chế tình trạng chatbot đưa ra thông tin vượt quá phạm vi knowledge base.

Ngoài ra, cần xây dựng bộ đánh giá hệ thống cho các kịch bản chính như tư vấn theo cửa hàng, so sánh sản phẩm, tham chiếu giá, tạo yêu cầu mua hàng, hỏi ngoài dữ liệu và lỗi AI Service. Các tiêu chí đánh giá nên bao gồm độ đúng của phạm vi dữ liệu, khả năng duy trì ngữ cảnh, khả năng hỏi làm rõ khi thiếu thông tin và thời gian phản hồi.

Cuối cùng, hệ thống có thể mở rộng thêm các chức năng vận hành thực tế như dashboard thống kê hội thoại, lead và yêu cầu mua hàng; backup dữ liệu; giám sát lỗi AI Service; cấu hình domain, HTTPS và bảo mật secret khi triển khai production. Về nghiệp vụ bán hàng, hệ thống có thể bổ sung quản lý giỏ hàng, đơn hàng, tồn kho, thanh toán hoặc tích hợp CRM.

\end{document}